\chapter{Avleggaren}

\begin{motto}
Frenologien forsvann ikkje fordi han var vond.\\ Han forsvann fordi nokon bygde nevrovitskapen ved sida av.
\end{motto}

Det finst ein kategori i vitskapshistoria som designfaget høyrer heime i, og som det aldri har vedkjent seg. Det er kategorien av felt som ein gong var respekterte, som hadde utøvarar, institusjonar, litteratur og prestisje, og som forsvann — ikkje fordi dei vart motbeviste i ein einaste avgjerande test, men fordi det vaks fram eit fag med eit fundament ved sida av dei, og det nye faget gjorde det gamle overflødig. Frenologien hadde tidsskrift, foreiningar, framståande utøvarar; han forsvann då nevrovitskapen kom. Alkemien hadde tusen års akkumulert praksis; han forsvann då kjemien kom. Vitalismen forklarte livet med ei eigen livskraft; han forsvann då biokjemien viste at livet er kjemi heile vegen ned. Astrologien kartla himmelen med stor presisjon; astronomien overtok himmelen og lét astrologien sitje att med teikna.

Det desse faga deler er ikkje at dei var dumme. Mange kloke menneske arbeidde i dei, og dei produserte reelle observasjonar. Det dei deler er ein \emph{struktur}: dei hadde ein katalog utan ein teori, eit mønstergjenkjennande blikk utan eit fundament som kunne grunngje mønstera, og difor kunne dei ikkje skilje dei sanne mønstera frå dei innbilte. Frenologen såg verkeleg samanhengar mellom hovud og karakter — han hadde berre ingen teori som kunne seie kva for samanhengar som var reelle. Det er nett designfaget si stode. Det ser verkelege mønster i form; det manglar berre teorien som kan seie kvifor dei er der, og difor kan det ikkje skilje innsikt frå smak, kvalitet frå mote, kompetanse frå konvensjon.

\section{Det før-paradigmatiske faget}

Thomas Kuhn gav oss ordet for tilstanden\kref{15}. Eit \emph{før-paradigmatisk} fag er eit fag før det har funne sitt fyrste delte rammeverk. Det kjenneteiknast av konkurrerande skular som ikkje kan avgjere usemjene sine, av at kvar generasjon byrjar grunnlaget på nytt, av at det ikkje finst eit delt sett problem alle arbeider på, og av at framgang ikkje akkumulerer. Kvar setning i denne skildringa passar på designfaget. Det har konkurrerande skular — Ulm mot Bauhaus mot den kritiske mot den spekulative mot design thinking — som ikkje kan avgjere usemjene sine, fordi det ikkje finst eit grunnlag å avgjere dei på. Kvar generasjon byggjer blikket sitt på nytt. Det finst ikkje eit delt sett problem; kvar forskar definerer sitt eige. Framgang akkumulerer ikkje; ho driv.

Men her er skilnaden som gjer designfaget sitt tilfelle alvorlegare enn eit vanleg ungt fag. Eit før-paradigmatisk fag er normalt eit fag på veg mot eit paradigme — kjemien før Lavoisier, biologien før Darwin, geologien før platetektonikken. Det er ein ungdomstilstand ein veks ut av. Designfaget er ikkje ungt. Det har vore akademisk i over hundre år. Det har hatt all den tid eit fag treng for å finne fundamentet sitt, og det har ikkje funne det — ikkje fordi fundamentet er uvanleg vanskeleg, men fordi faget har bygd, som dei føregåande kapitla viste, eit heilt apparat av sjølvforsvar som gjer at det ikkje treng å leite. Det har gjort den før-paradigmatiske tilstanden permanent ved å omdøype han til faget sin natur. «Design er ein eigen kunnskapsform.» «Designproblem er vrange.» Det er ikkje skildringar av eit fag på veg mot eit paradigme. Det er erklæringar om at faget aldri har tenkt å finne eitt. Eit før-paradigmatisk fag som nektar å bli paradigmatisk, er ikkje lenger ungt. Det er avleggar i vente.

\section{Kvifor avviklinga kjem no}

Avleggarane forsvann ikkje når dei vart motbeviste; dei forsvann når verda fekk bruk for noko dei ikkje kunne levere, og eit anna fag kunne. Den krafta er no på veg mot designfaget, og ho kjem frå materialet sjølv.

I heile faget si historie har formgjevaren arbeidd med passivt stoff: tre, metall, plast, pikslar — materiale som gjorde nett det dei vart sagt, og ikkje meir. For slikt stoff kunne eit fag med berre intuisjon og handverk så vidt halde det gåande, fordi konsekvensane av å mangle ein teori var avgrensa til det stoffet ikkje sa imot. No endrar materialet seg under hendene våre. Vi byggjer med agentisk materiale: celler som navigerer, levande system som reorganiserer seg, sjølvmonterande stoff\kref{14}. Vi byggjer med generative modellar som navigerer latente rom på eiga hand. Vi byggjer med system der celle, materiale, verktøy, marknad og modell alle er agentar med eigne horisontar som deformerer landskapet for kvarandre samstundes. For å formgje med slikt materiale treng du presis det faget ikkje har: ein modell av kva ein agent er, kva eit seleksjonstrykk er, kva eit formrom er, korleis kompetanse fordeler seg over skalaer. Du kan ikkje intuere deg fram med ein celle som har sin eigen agenda. Du kan ikkje halde kritikk over ein generativ modell. Det mønstergjenkjennande blikket, faget sitt einaste reiskap, er verdlaust andsynes materiale som svarar tilbake.

Det er dette som gjer samanlikninga med frenologien presis snarare enn polemisk. Frenologen kunne halde fram så lenge alt han hadde var skallar å kjenne på. Gjev honom ein hjerneskannar, og reiskapen avslører straks at teorien bak handa hans aldri forstod noko. Designfaget får no sine hjerneskannarar — verktøy så kraftige og materiale så tydeleg agentisk at fråvêret av ein teori ikkje lenger kan gøymast bak handverket. Og samstundes vert sjølve handverket, faget sin siste eigentlege kompetanse, automatisert: generative system gjer no det studentane brukte år på å lære i atelieret, raskare og i uendeleg mengd. Når maskina kan produsere formene, sit faget att med berre éin ting det kunne hatt og ikkje har: teorien om kvifor ei form er betre enn ei anna. Det er det einaste som ikkje let seg automatisere bort, og det er nett det faget aldri bygde.

\section{To måtar å bli avleggar på}

Lat meg vere presis om kva eg spår, for eg spår ikkje at design forsvinn. Verda vil alltid trenge at nokon gjev form til ting; det er like umisteleg som at verda treng at nokon lækjer sjuke. Det eg spår er at \emph{faget slik det er} — den før-paradigmatiske disiplinen som deler ut grader i ein kompetanse han ikkje kan definere — kjem til å bli avleggar, slik medisinen før fysiologien vart avleggar medan lækjinga heldt fram. Formgjevinga overlever. Spørsmålet er kva fag som kjem til å eige henne.

Og her er det to vegar, og berre to. Anten finn formgjevinga sitt fundament innanfrå — designfaget gjer endeleg den overgangen det har utsett i hundre år, byggjer eller adopterer ein presis teori om form, og vert eit paradigmatisk fag som fortener plassen sin ved universitetet. Eller så vert formgjevinga overteken utanfrå — av informatikken, av materialvitskapen, av eit komande felt for agentisk formgjeving som har det fundamentet designfaget mangla — og designfaget sit att med teikna, slik astrologen sat att med stjerneteikna då astronomen overtok himmelen. Begge vegar endar den noverande disiplinen. Den eine let han bli til noko betre. Den andre let han bli til ein fotnote. Det faget ikkje lenger har, er den tredje vegen: å halde fram som før. Den vegen stengjer materialet og maskinene no, samstundes, og utan å spørje om lov.

\vignett

Klagemålet er ført. Faget kan ikkje grunngje dommane sine, ikkje fordøye låna sine, ikkje gjere forskinga kumulativ, ikkje seie nei, og ikkje sjå seg sjølv klart; og historia har ein hard plass reservert for felt i den tilstanden. Det står att eitt spørsmål, og berre eitt: finst det ein grunn å byggje på, om faget vil? Det gjer det. Det siste kapitlet peikar mot han.
